% ===========================================================================
%  Chapitre 16 — Géométrie dans l'espace
%  PE 2026, composante GÉOMÉTRIE
% ===========================================================================
\bmtchapitre{Géométrie dans l'espace}
  {Résoudre un problème dans l'espace : décomposition de vecteurs, produits
   scalaire et vectoriel, équations de droites et de plans, positions
   relatives.}
  {Géométrie \textbullet\ environ 10 heures}

Passer du plan à l'espace ajoute une coordonnée, et cela change moins de
choses qu'on ne le craint : les vecteurs se manipulent de la même façon, le
produit scalaire garde la même définition, les équations restent linéaires. Ce
qui change, c'est le vocabulaire des positions relatives — deux droites de
l'espace peuvent n'être ni sécantes ni parallèles — et l'apparition d'un outil
nouveau, le produit vectoriel.

Ce produit vectoriel est la clé du chapitre. Il fabrique, à partir de deux
vecteurs, un troisième qui leur est orthogonal ; il donne donc directement le
vecteur normal d'un plan défini par trois points. L'exercice de géométrie dans
l'espace du baccalauréat, deux points sur vingt, commence presque toujours par
ce calcul.

% ---------------------------------------------------------------------------
\begin{revision}
Géométrie plane et calcul vectoriel de première.

\begin{enumerate}[label=\textbf{\arabic*.}, leftmargin=*, itemsep=0.35em]
  \item Rappeler l'expression analytique du produit scalaire dans le plan.
  \item Que signifie $\scal{\vecu}{\vecv} = 0$ ?
  \item Calculer la norme du vecteur de coordonnées $(2\,;-1\,;2)$.
  \item Rappeler la relation de Chasles et la définition de deux vecteurs
        colinéaires.
  \item Écrire l'équation cartésienne d'une droite du plan de vecteur normal
        $\vecn(a\,;b)$ passant par $A$.
\end{enumerate}
\end{revision}

\begin{automatismes}
Sans calculatrice, sans brouillon.

\begin{enumerate}[label=\textbf{\alph*)}, leftmargin=*, itemsep=0.25em]
  \item $\scal{(1;2;3)}{(1;0;-1)}$
  \item Norme de $(3\,;0\,;4)$
  \item $(1;0;0) \wedge (0;1;0)$
  \item Vecteur normal du plan $2x-y+5z-3 = 0$
  \item Le point $(1\,;1\,;0)$ appartient-il à $x+y-z = 2$ ?
  \item Deux vecteurs colinéaires : leur produit vectoriel
  \item Direction de la droite $x = 1+t$, $y = -t$, $z = 2t$
  \item Distance de l'origine au plan $z = 3$
\end{enumerate}
\end{automatismes}

\begin{corrige}[automatismes]
a) $-2$ \quad b) $5$ \quad c) $(0;0;1)$ \quad d) $(2;-1;5)$ \quad
e) oui \quad f) le vecteur nul \quad g) $(1;-1;2)$ \quad h) $3$.
\end{corrige}

% ===========================================================================
\section{Vecteurs de l'espace et décomposition}

\begin{definition}[vecteurs coplanaires, base]
Trois vecteurs de l'espace sont \textbf{coplanaires} s'il existe un plan dont
ils dirigent des droites. Trois vecteurs non coplanaires forment une
\textbf{base} de l'espace : tout vecteur $\vect{w}$ s'écrit alors de manière
unique
\[
  \vect{w} = x\,\veci+y\,\vecj+z\,\veck .
\]
\end{definition}

\begin{propriete}[repère orthonormé]
Un repère $\left(O\,;\veci,\vecj,\veck\right)$ est \textbf{orthonormé} si les
trois vecteurs sont deux à deux orthogonaux et de norme $1$. Dans un tel
repère, pour $\vecu(x\,;y\,;z)$,
\[
  \norme{\vecu} = \sqrt{x^{2}+y^{2}+z^{2}},
  \qquad
  AB = \sqrt{(x_{B}-x_{A})^{2}+(y_{B}-y_{A})^{2}+(z_{B}-z_{A})^{2}} .
\]
\end{propriete}

\begin{remarque}
Trois points de l'espace définissent un plan si et seulement s'ils ne sont pas
alignés, c'est-à-dire si $\vect{AB}$ et $\vect{AC}$ ne sont pas colinéaires.
C'est la toute première vérification de tout exercice d'espace, et elle se
fait par le produit vectoriel.
\end{remarque}

\begin{visualisation}[le repère de l'espace]
\begin{center}
\begin{tikzpicture}[x={(-0.5cm,-0.28cm)}, y={(1cm,0cm)}, z={(0cm,1cm)},
                     font=\footnotesize]
  \draw[bmtGris, opacity=0.5] (0,0,0) -- (2.4,0,0) -- (2.4,2.6,0) -- (0,2.6,0) -- cycle;
  \draw[bmtGris, -{Stealth[length=2.6mm]}] (0,0,0) -- (2.9,0,0) node[anchor=north east]{$\veci$};
  \draw[bmtGris, -{Stealth[length=2.6mm]}] (0,0,0) -- (0,3.1,0) node[anchor=north west]{$\vecj$};
  \draw[bmtGris, -{Stealth[length=2.6mm]}] (0,0,0) -- (0,0,2.7) node[anchor=south]{$\veck$};
  \node[bmtGris, anchor=east] at (0.15,-0.15,0) {$O$};
  \draw[bmtGris, dotted] (1.6,0,0) -- (1.6,1.8,0) -- (0,1.8,0);
  \draw[bmtGris, dotted] (1.6,1.8,0) -- (1.6,1.8,1.4);
  \node[bmtGris, anchor=north] at (1.6,0,0) {$x$};
  \node[bmtGris, anchor=north west] at (0,1.8,0) {$y$};
  \node[bmtGris, anchor=east] at (0,0,1.4) {$z$};
  \draw[bmtGris, dotted] (0,0,1.4) -- (1.6,1.8,1.4);
  \filldraw[bmtIndigo] (1.6,1.8,1.4) circle (2pt) node[anchor=south west] {$M(x\,;y\,;z)$};
\end{tikzpicture}
\end{center}
Dans le repère $\left(O\,;\veci,\vecj,\veck\right)$, chaque point $M$ se lit
par ses trois projections sur les axes. Le parallélépipède en pointillés est
le repère de lecture habituel : on avance de $x$ selon $\veci$, de $y$ selon
$\vecj$, puis de $z$ selon $\veck$.
\end{visualisation}

% ===========================================================================
\section{Produit scalaire}

\begin{definition}[produit scalaire dans l'espace]
Pour $\vecu(x\,;y\,;z)$ et $\vecv(x'\,;y'\,;z')$ dans un repère orthonormé,
\[
  \scal{\vecu}{\vecv} = xx'+yy'+zz'
  = \norme{\vecu}\,\norme{\vecv}\cos\theta,
\]
où $\theta$ est l'angle géométrique entre les deux vecteurs.
\end{definition}

\begin{propriete}[orthogonalité]
Deux vecteurs non nuls sont orthogonaux si et seulement si leur produit
scalaire est nul. Une droite est orthogonale à un plan si et seulement si son
vecteur directeur est orthogonal à deux vecteurs non colinéaires du plan.
\end{propriete}

\begin{visualisation}[l'angle entre deux vecteurs]
\begin{center}
\begin{tikzpicture}[x=1cm, y=1cm, font=\footnotesize]
  \draw[bmtIndigo, line width=1.1pt, -{Stealth[length=2.6mm]}]
    (0,0) -- (3.2,0) node[anchor=north west]{$\vecu$};
  \draw[bmtLaterite, line width=1.1pt, -{Stealth[length=2.6mm]}]
    (0,0) -- (35:2.6) node[anchor=south west]{$\vecv$};
  \draw[bmtGris, dashed] (35:2.6) -- ({2.6*cos(35)},0);
  \draw[bmtVetiver, line width=2.2pt] (0,0) -- ({2.6*cos(35)},0);
  \node[bmtVetiver, anchor=north] at ({0.5*2.6*cos(35)},-0.05) {projeté};
  \draw[bmtGris] (0.55,0) arc (0:35:0.55);
  \node[bmtGris, anchor=west] at (17:0.75) {$\theta$};
\end{tikzpicture}
\end{center}
Le produit scalaire $\scal{\vecu}{\vecv}$ mesure, à la norme de $\vecu$
près, la longueur algébrique du projeté de $\vecv$ sur $\vecu$ : positif si
$\theta$ est aigu, négatif s'il est obtus, nul à angle droit. C'est cette
lecture qui rend le critère d'orthogonalité évident.
\end{visualisation}

% ===========================================================================
\section{Produit vectoriel}

\begin{definition}[produit vectoriel]
Le \textbf{produit vectoriel} de $\vecu(x\,;y\,;z)$ et $\vecv(x'\,;y'\,;z')$
est le vecteur
\[
  \pvect{\vecu}{\vecv} =
  \begin{pmatrix}
    yz'-zy' \\ zx'-xz' \\ xy'-yx'
  \end{pmatrix}.
\]
\end{definition}

\begin{propriete}[caractérisation géométrique]
$\pvect{\vecu}{\vecv}$ est orthogonal à $\vecu$ et à $\vecv$ ; sa norme vaut
$\norme{\vecu}\,\norme{\vecv}\,\abs{\sin\theta}$, qui est l'aire du
parallélogramme construit sur les deux vecteurs. Il est nul si et seulement si
$\vecu$ et $\vecv$ sont colinéaires. Enfin
$\pvect{\vecv}{\vecu} = -\pvect{\vecu}{\vecv}$ : le produit vectoriel n'est
pas commutatif.
\end{propriete}

\begin{visualisation}[le vecteur qui sort du plan]
\begin{center}
\begin{tikzpicture}[x=1cm, y=1cm, font=\footnotesize]
  \fill[bmtVetiver, opacity=0.12] (-0.6,-0.4) -- (4.6,-0.4) -- (5.6,1.3)
    -- (0.4,1.3) -- cycle;
  \draw[bmtVetiver, opacity=0.5] (-0.6,-0.4) -- (4.6,-0.4) -- (5.6,1.3)
    -- (0.4,1.3) -- cycle;
  \coordinate (O) at (1.2,0.1);
  \draw[bmtIndigo, line width=1.1pt, -{Stealth[length=2.6mm]}]
    (O) -- ++(2.6,0.25) node[anchor=north west] {$\vecu$};
  \draw[bmtLaterite, line width=1.1pt, -{Stealth[length=2.6mm]}]
    (O) -- ++(1.05,0.95) node[anchor=south east] {$\vecv$};
  \draw[bmtEncre, line width=1.2pt, -{Stealth[length=2.8mm]}]
    (O) -- ++(0,2.35) node[anchor=south] {$\pvect{\vecu}{\vecv}$};
  \draw[bmtGris, dashed] ($(O)+(2.6,0.25)$) -- ($(O)+(3.65,1.2)$)
    -- ($(O)+(1.05,0.95)$);
  \node[bmtGris] at ($(O)+(1.85,0.62)$) {aire};
\end{tikzpicture}
\end{center}

Le produit vectoriel sort du plan des deux vecteurs : c'est précisément ce qui
en fait l'outil du chapitre. Sa direction donne le vecteur normal d'un plan,
sa norme mesure l'aire du parallélogramme dessiné en pointillés, et sa
nullité signale que les deux vecteurs sont colinéaires — donc qu'ils ne
définissent aucun plan.
\end{visualisation}

\begin{methode}{calculer un produit vectoriel sans se tromper}
On écrit les deux triplets l'un sous l'autre, puis on répète les deux
premières lignes en dessous. Chaque coordonnée du résultat s'obtient en
« croisant » les deux lignes qui l'encadrent : produit descendant moins
produit montant. La vérification est immédiate : le produit scalaire du
résultat avec chacun des deux vecteurs doit être nul.
\end{methode}

\begin{exemple}[un plan par trois points]
Soient $A(2\,;3\,;1)$, $B(-1\,;1\,;1)$ et $C(-3\,;1\,;0)$. Alors
$\vect{AB}(-3\,;-2\,;0)$ et $\vect{AC}(-5\,;-2\,;-1)$, et
\[
  \pvect{\vect{AB}}{\vect{AC}} =
  \begin{pmatrix}
    (-2)(-1)-0 \times(-2) \\ 0 \times(-5)-(-3)(-1) \\
    (-3)(-2)-(-2)(-5)
  \end{pmatrix}
  = \begin{pmatrix} 2 \\ -3 \\ -4 \end{pmatrix}.
\]
Ce vecteur n'est pas nul : les trois points ne sont pas alignés et
définissent bien un plan, d'équation $2x-3y-4z+d = 0$. Le point $A$ donne
$4-9-4+d = 0$, soit $d = 9$ :
\[
  (ABC) : \quad 2x-3y-4z+9 = 0 .
\]
\emph{Vérification.} $B$ : $-2-3-4+9 = 0$. $C$ : $-6-3-0+9 = 0$.
\end{exemple}

\begin{entrainement}[produits scalaire et vectoriel]
\exo[\niveau{1}] Calculer $\scal{\vecu}{\vecv}$ et $\pvect{\vecu}{\vecv}$ pour
$\vecu(1\,;2\,;-1)$ et $\vecv(3\,;0\,;2)$.

\exo[\niveau{1}] Les vecteurs $(2\,;-4\,;6)$ et $(1\,;-2\,;3)$ sont-ils
colinéaires ?

\exo[\niveau{2}] Calculer l'aire du triangle $ABC$ pour $A(1\,;0\,;0)$,
$B(0\,;2\,;0)$ et $C(0\,;0\,;3)$.

\exo[\niveau{2}] Montrer que $\pvect{\vecu}{\vecu}$ est le vecteur nul.

\exo[\niveau{3}] Montrer que $\norme{\pvect{\vecu}{\vecv}}^{2}
+\left(\scal{\vecu}{\vecv}\right)^{2}
= \norme{\vecu}^{2}\norme{\vecv}^{2}$.
\end{entrainement}

\begin{corrige}
\textbf{1.} $\scal{\vecu}{\vecv} = 3+0-2 = 1$ et
$\pvect{\vecu}{\vecv} = \left(4\,;-5\,;-6\right)$. On vérifie que ce dernier
est orthogonal aux deux : $4-10+6 = 0$ et $12+0-12 = 0$.

\textbf{2.} $(2\,;-4\,;6) = 2\times(1\,;-2\,;3)$ : les deux vecteurs sont
colinéaires, et leur produit vectoriel est nul.

\textbf{3.} $\vect{AB}(-1\,;2\,;0)$ et $\vect{AC}(-1\,;0\,;3)$, donc
$\pvect{\vect{AB}}{\vect{AC}} = (6\,;3\,;2)$, de norme
$\sqrt{36+9+4} = 7$. L'aire du triangle est la moitié de cette norme, soit
$\dfrac72$.

\textbf{4.} Le produit vectoriel de deux vecteurs colinéaires est nul, et un
vecteur est toujours colinéaire à lui-même. On peut aussi le lire sur les
coordonnées : chaque composante est une différence de deux produits égaux.

\textbf{5.} En notant $\theta$ l'angle des deux vecteurs,
$\norme{\pvect{\vecu}{\vecv}} = \norme{\vecu}\norme{\vecv}\abs{\sin\theta}$ et
$\scal{\vecu}{\vecv} = \norme{\vecu}\norme{\vecv}\cos\theta$. La somme des
carrés vaut donc
$\norme{\vecu}^{2}\norme{\vecv}^{2}\left(\sin^{2}\theta+\cos^{2}\theta\right)
= \norme{\vecu}^{2}\norme{\vecv}^{2}$.
\end{corrige}

% ===========================================================================
\section{Équations de droites et de plans}

\begin{propriete}[équation cartésienne d'un plan]
Un plan de vecteur normal $\vecn(a\,;b\,;c)$ non nul a une équation de la
forme
\[
  ax+by+cz+d = 0,
\]
et réciproquement toute équation de ce type définit un plan de vecteur normal
$\vecn(a\,;b\,;c)$.
\end{propriete}

\begin{propriete}[représentation paramétrique d'une droite]
La droite passant par $A(x_{0}\,;y_{0}\,;z_{0})$ et dirigée par
$\vecu(\alpha\,;\beta\,;\gamma)$ est l'ensemble des points de coordonnées
\[
  \begin{cases}
    x = x_{0}+\alpha t \\
    y = y_{0}+\beta t \\
    z = z_{0}+\gamma t
  \end{cases}
  \qquad t \in \R .
\]
\end{propriete}

\begin{visualisation}[ce qui engendre un plan, ce qui engendre une droite]
\begin{center}
\begin{tikzpicture}[x=1cm, y=1cm, font=\footnotesize]
  \begin{scope}
    \fill[bmtVetiver, opacity=0.12] (-0.6,-0.4) -- (3.4,-0.4) -- (4.2,1.1)
      -- (0.2,1.1) -- cycle;
    \draw[bmtVetiver, opacity=0.6] (-0.6,-0.4) -- (3.4,-0.4) -- (4.2,1.1)
      -- (0.2,1.1) -- cycle;
    \node[bmtVetiver, anchor=west] at (3.6,-0.15) {$(P)$};
    \coordinate (A) at (1.6,0.25);
    \filldraw[bmtEncre] (A) circle (2pt) node[anchor=north east]{$A$};
    \draw[bmtLaterite, line width=1.2pt, -{Stealth[length=2.8mm]}]
      (A) -- ++(0.1,1.7) node[anchor=south]{$\vecn(a\,;b\,;c)$};
    \node[bmtGris, anchor=north, font=\footnotesize] at (1.8,-1.1)
      {$ax+by+cz+d=0$};
  \end{scope}
  \begin{scope}[xshift=6.4cm]
    \draw[bmtGris, line width=0.6pt] (-0.8,-0.9) -- (2.6,1.6);
    \coordinate (A2) at (0.3,-0.15);
    \coordinate (Mt) at (1.7,0.85);
    \filldraw[bmtEncre] (A2) circle (2pt) node[anchor=north]{$A(x_{0};y_{0};z_{0})$};
    \filldraw[bmtIndigo] (Mt) circle (1.8pt) node[anchor=south west]{$M$};
    \draw[bmtLaterite, line width=1.2pt, -{Stealth[length=2.8mm]}]
      (A2) -- ++(0.9,0.6) node[anchor=south east]{$\vecu$};
    \node[bmtGris, anchor=north, font=\footnotesize] at (1,-1.1)
      {$\vect{AM}=t\,\vecu$};
  \end{scope}
\end{tikzpicture}
\end{center}
Un plan est engendré par un point et un vecteur normal : l'équation
$ax+by+cz+d=0$ n'est jamais qu'une reformulation de
$\scal{\vect{AM}}{\vecn} = 0$. Une droite est engendrée par un point et un
vecteur directeur : tout point $M$ de la droite s'obtient en avançant de $t$
fois $\vecu$ depuis $A$ — c'est exactement le système paramétrique.
\end{visualisation}

\begin{propriete}[distance d'un point à un plan]
La distance du point $A(x_{A}\,;y_{A}\,;z_{A})$ au plan
$(P) : ax+by+cz+d = 0$ vaut
\[
  d\left(A,(P)\right)
  = \frac{\abs{ax_{A}+by_{A}+cz_{A}+d}}{\sqrt{a^{2}+b^{2}+c^{2}}} .
\]
\end{propriete}

\begin{entrainement}[équations]
\exo[\niveau{1}] Donner une équation du plan passant par $A(1\,;0\,;2)$ et de
vecteur normal $\vecn(2\,;-1\,;3)$.

\exo[\niveau{2}] Donner une représentation paramétrique de la droite passant
par $A(1\,;-1\,;1)$ et $B(3\,;2\,;1)$.

\exo[\niveau{2}] Calculer la distance du point $A(2\,;2\,;1)$ au plan
$3x+2y+6z+33 = 0$.

\exo[\niveau{3}] Déterminer une équation du plan contenant la droite
$x = t$, $y = 1+t$, $z = 2-t$ et passant par l'origine.
\end{entrainement}

\begin{corrige}
\textbf{6.} $2(x-1)-(y-0)+3(z-2) = 0$, soit $2x-y+3z-8 = 0$.

\textbf{7.} $\vect{AB}(2\,;3\,;0)$ est vecteur directeur, d'où
$x = 1+2t$, $y = -1+3t$, $z = 1$, avec $t \in \R$. La cote est constante : la
droite est horizontale, contenue dans le plan $z = 1$.

\textbf{8.} $d = \dfrac{\abs{3\times2+2\times2+6\times1+33}}{\sqrt{9+4+36}}
= \dfrac{49}{7} = 7$.

\textbf{9.} La droite passe par $P(0\,;1\,;2)$ et a pour vecteur directeur
$\vect{d}(1\,;1\,;-1)$. Le plan cherché contient $O$ et $P$, donc le vecteur
$\vect{OP}(0\,;1\,;2)$, ainsi que $\vect{d}$ : un vecteur normal est
$\pvect{\vect{d}}{\vect{OP}} = (3\,;-2\,;1)$. Passant par l'origine, le plan a
pour équation $3x-2y+z = 0$ ; on vérifie que $P$ y appartient et que $\vect{d}$
est bien orthogonal au vecteur normal.
\end{corrige}

% ===========================================================================
\section{Positions relatives}

\begin{propriete}[deux plans]
Deux plans de vecteurs normaux $\vecn$ et $\vecn'$ sont parallèles si et
seulement si $\vecn$ et $\vecn'$ sont colinéaires. Ils sont alors confondus ou
strictement parallèles ; sinon ils se coupent suivant une droite.
\end{propriete}

\begin{propriete}[un plan et une droite]
Soit $(P)$ de vecteur normal $\vecn$ et $(D)$ de vecteur directeur $\vecu$.
\begin{itemize}[leftmargin=*, itemsep=0.1em]
  \item Si $\scal{\vecn}{\vecu} \neq 0$ : la droite coupe le plan en un point
        unique.
  \item Si $\scal{\vecn}{\vecu} = 0$ : la droite est parallèle au plan —
        contenue dedans si l'un de ses points y appartient, strictement
        parallèle sinon.
  \item Si $\vecu$ et $\vecn$ sont colinéaires, la droite est perpendiculaire
        au plan.
\end{itemize}
\end{propriete}

\begin{propriete}[deux droites de l'espace]
Deux droites de l'espace sont parallèles si leurs vecteurs directeurs sont
colinéaires. Sinon, elles sont sécantes ou \textbf{non coplanaires} : elles
sont coplanaires si et seulement si les trois vecteurs $\vecu$, $\vecu'$ et
$\vect{AA'}$ — où $A$ et $A'$ sont un point de chacune — sont coplanaires,
c'est-à-dire si
$\scal{\left(\pvect{\vecu}{\vecu'}\right)}{\vect{AA'}} = 0$.
\end{propriete}

\begin{avertissement}
Dans l'espace, « non parallèles » n'entraîne pas « sécantes ». Deux droites
peuvent n'avoir aucun point commun sans être parallèles : on dit alors
qu'elles sont non coplanaires. C'est la situation la plus fréquemment
proposée au baccalauréat, et l'oubli de ce cas fausse toute la question.
\end{avertissement}

\begin{visualisation}[deux droites qui ne se rencontrent jamais]
\begin{center}
\begin{tikzpicture}[x={(-0.5cm,-0.28cm)}, y={(1cm,0cm)}, z={(0cm,1cm)},
                     font=\footnotesize]
  \fill[bmtGris, opacity=0.10] (-0.6,-0.6,0) -- (3.2,-0.6,0)
    -- (3.2,2.8,0) -- (-0.6,2.8,0) -- cycle;
  \draw[bmtGris, opacity=0.45] (-0.6,-0.6,0) -- (3.2,-0.6,0)
    -- (3.2,2.8,0) -- (-0.6,2.8,0) -- cycle;
  \draw[bmtIndigo, line width=1.2pt] (0.2,0.3,0) -- (2.8,1.8,0)
    node[anchor=west]{$(D_{1})$};
  \draw[bmtLaterite, line width=1.2pt] (0.3,2.0,1.7) -- (2.6,0.2,2.1)
    node[anchor=south]{$(D_{2})$};
  \draw[bmtGris, dotted] (0.3,2.0,1.7) -- (0.3,2.0,0);
  \draw[bmtGris, dotted] (2.6,0.2,2.1) -- (2.6,0.2,0);
  \filldraw[bmtGris] (0.3,2.0,0) circle (1.3pt);
  \filldraw[bmtGris] (2.6,0.2,0) circle (1.3pt);
\end{tikzpicture}
\end{center}
$(D_{1})$ repose sur le plan de fond, $(D_{2})$ flotte au-dessus, dans une
autre direction. Leurs projections verticales, en pointillés, ne se
confondent avec $(D_{1})$ qu'aux deux pieds marqués : les deux droites ne
sont ni parallèles ni sécantes, elles sont \textbf{non coplanaires}. C'est
la situation la plus fréquente de l'exercice de baccalauréat, et la plus
souvent oubliée.
\end{visualisation}

\begin{methode}{montrer que deux droites sont non coplanaires}
\begin{enumerate}[label=\textbf{\arabic*.}, leftmargin=*, itemsep=0.15em]
  \item Vérifier que les vecteurs directeurs ne sont pas colinéaires : les
        droites ne sont donc pas parallèles.
  \item Choisir un point $A$ sur la première, $A'$ sur la seconde.
  \item Calculer $\scal{\left(\pvect{\vecu}{\vecu'}\right)}{\vect{AA'}}$. S'il
        est non nul, les droites ne sont pas coplanaires ; s'il est nul, elles
        sont sécantes.
\end{enumerate}
\end{methode}

\begin{exemple}[deux droites non coplanaires]
Soient $(D_{1}) : x = 1+t$, $y = -1-2t$, $z = 1+3t$ et
$(D_{2}) : x = 3+2m$, $y = m$, $z = 1-4m$.

Les directions $\vecu(1\,;-2\,;3)$ et $\vecu'(2\,;1\,;-4)$ ne sont pas
colinéaires. On prend $A(1\,;-1\,;1)$ et $A'(3\,;0\,;1)$, donc
$\vect{AA'}(2\,;1\,;0)$. Alors
\[
  \pvect{\vecu}{\vecu'} = (5\,;10\,;5),
  \qquad
  \scal{(5;10;5)}{(2;1;0)} = 10+10+0 = 20 \neq 0 .
\]
Les deux droites ne sont donc pas coplanaires.
\end{exemple}

% ===========================================================================
\begin{annales}
L'exercice de géométrie dans l'espace vaut deux ou trois points et se traite
en vingt minutes lorsqu'on maîtrise le produit vectoriel. Les énoncés qui
suivent sont repris de sujets réellement posés ; leur provenance est indiquée
à chaque fois.

\exoann{Baccalauréat, Madagascar, série S, session 2021 — exercice 2}
L'espace est rapporté à un repère orthonormé direct. On considère
$A(1\,;-1\,;1)$, $B(3\,;2\,;1)$ et $C(2\,;-1\,;0)$.
\begin{enumerate}[label=\textbf{\arabic*)}, leftmargin=*, itemsep=0.15em]
  \item Vérifier que l'équation du plan $(ABC)$ est $3x-2y+3z-8 = 0$.
  \item Écrire une équation cartésienne du plan $(P)$ contenant $B$ et de
        vecteur normal $\vect{AC}$.
\end{enumerate}

\exoann{Baccalauréat, Madagascar, série S, session 2021 — exercice 2, suite}
\begin{enumerate}[label=\textbf{\alph*)}, leftmargin=*, itemsep=0.15em]
  \item Démontrer que l'intersection des plans $(ABC)$ et $(P)$ est la droite
        $(D)$ de représentation paramétrique
        $x = t+2$, $y = 3t-1$, $z = t$, avec $t \in \R$.
  \item Montrer que $C$ est le projeté orthogonal de $A$ sur $(D)$ et en
        déduire la distance du point $A$ à la droite $(D)$.
\end{enumerate}

\exoann{Baccalauréat, Madagascar, série S, session 2021 — exercice 2, fin}
Soit $(Q)$ le plan d'équation $x-y+z = 1$. Déterminer les coordonnées du point
$I$, intersection des trois plans $(ABC)$, $(P)$ et $(Q)$.

\exoann{Baccalauréat, Madagascar, série S, session 2022 — exercice 2}
L'espace est rapporté à un repère orthonormé direct. On donne
$A(2\,;3\,;1)$, $B(-1\,;1\,;1)$, $C(-3\,;1\,;0)$ et $D(4\,;-1\,;3)$.
\begin{enumerate}[label=\textbf{\arabic*)}, leftmargin=*, itemsep=0.15em]
  \item Calculer $\pvect{\vect{AB}}{\vect{AC}}$.
  \item Écrire une équation cartésienne du plan $(ABC)$.
  \item Donner une représentation paramétrique de la droite $(\Delta)$ de
        vecteur directeur $\vecu(1\,;2\,;-1)$ passant par $D$.
  \item Montrer que le plan $(ABC)$ et la droite $(\Delta)$ sont parallèles.
\end{enumerate}

\exoann{Baccalauréat, Madagascar, série S, session 2023 — exercice 2}
L'espace est muni d'un repère orthonormé. On donne le point $A(2\,;2\,;1)$ et
le plan $(\mathcal P)$ d'équation $3x+2y+6z+33 = 0$.
\begin{enumerate}[label=\textbf{\arabic*)}, leftmargin=*, itemsep=0.15em]
  \item Vérifier que $(\mathcal P)$ est strictement parallèle à un plan
        $(\mathcal Q)$ contenant $A$ et de vecteur normal
        $\vecn(-3\,;-2\,;-6)$.
  \item Déterminer la distance qui sépare ces deux plans.
  \item Soit $H(-1\,;0\,;-5)$. Justifier que $H \in (\mathcal P)$ et que
        $(AH)$ et $(\mathcal P)$ sont orthogonaux, puis calculer
        $\norme{\vect{AH}}$.
\end{enumerate}

\exoann{Baccalauréat, Madagascar, série S, session 2026 — exercice 2}
L'espace orienté est muni d'un repère orthonormé direct. Le plan $(P)$
contient $A(1\,;-2\,;2)$ et est dirigé par $\vecu(-1\,;2\,;1)$ et
$\vecv(-4\,;2\,;-2)$. La droite $(D)$ a pour représentation paramétrique
$x = -2t+1$, $y = -2t+3$, $z = 2t-5$.
\begin{enumerate}[label=\textbf{\arabic*)}, leftmargin=*, itemsep=0.15em]
  \item Montrer que $(P)$ a pour équation cartésienne $x+y-z+3 = 0$.
  \item Vérifier que $(P)$ et $(D)$ sont perpendiculaires.
  \item Déterminer les coordonnées du point $F$, intersection de $(P)$ et
        $(D)$.
\end{enumerate}

\exoann{Baccalauréat blanc, lycée Philibert Tsiranana, Terminale S, 2021-2022 — exercice 2}
Dans un repère orthonormé direct, on considère
\[
  (D_{1}) : \begin{cases} x = 1+t \\ y = -1-2t \\ z = 1+3t \end{cases}
  \qquad\text{et}\qquad
  (D_{2}) : \begin{cases} x = 3+2m \\ y = m \\ z = 1-4m \end{cases}
\]
Montrer que $(D_{1})$ et $(D_{2})$ ne sont pas coplanaires.

\exoann{Baccalauréat blanc, lycée Philibert Tsiranana, Terminale S, 2021-2022 — exercice 2, suite}
Soit $(P)$ le plan passant par $A(1\,;-1\,;1)$ et dirigé par
$\vecu_{1}(1\,;-2\,;3)$ et $\vecu_{2}(2\,;1\,;-4)$.
\begin{enumerate}[label=\textbf{\alph*)}, leftmargin=*, itemsep=0.15em]
  \item Déterminer les composantes de $\vecn = \pvect{\vecu_{1}}{\vecu_{2}}$.
  \item Écrire une équation cartésienne de $(P)$.
  \item Vérifier que $(D_{1})$ est une droite du plan $(P)$.
\end{enumerate}

\exoann{Baccalauréat, Cameroun, série C, session 2023 — exercice 3}
L'espace est muni d'un repère orthonormé direct. Soient $A(-1\,;-1\,;0)$,
$B(0\,;0\,;2)$ et $C(-1\,;1\,;2)$. Montrer que $A$, $B$ et $C$ définissent un
plan, puis en donner une équation cartésienne.

\exoann{Baccalauréat, Côte d'Ivoire, série C, session 2023 — exercice 1}
Dire si la proposition suivante est vraie ou fausse, en justifiant.
« Dans un repère orthonormé, pour le plan d'équation $2x+3y+4z-8 = 0$, le
vecteur $\vecn\!\left(1\,;\frac32\,;2\right)$ est un vecteur normal. »
\end{annales}

\begin{corrige}[annales]
\textbf{1.} $\vect{AB}(2\,;3\,;0)$ et $\vect{AC}(1\,;0\,;-1)$. Le vecteur
$(3\,;-2\,;3)$ vérifie $6-6+0 = 0$ et $3+0-3 = 0$ : il est normal au plan.
L'équation est donc $3x-2y+3z+d = 0$, et $A$ donne $3+2+3+d = 0$, soit
$d = -8$. Pour $(P)$, le vecteur normal est $\vect{AC}(1\,;0\,;-1)$ et
$B(3\,;2\,;1)$ donne $3-1+d = 0$ : $(P) : x-z-2 = 0$.

\textbf{2. a)} Un point de $(D)$ vérifie
$3(t+2)-2(3t-1)+3t-8 = 3t+6-6t+2+3t-8 = 0$ et
$(t+2)-t-2 = 0$ : il appartient aux deux plans. Comme les deux plans ne sont
pas parallèles — leurs normales $(3\,;-2\,;3)$ et $(1\,;0\,;-1)$ ne sont pas
colinéaires —, leur intersection est une droite, et c'est $(D)$.
\textbf{b)} $C(2\,;-1\,;0)$ correspond à $t = 0$, donc $C \in (D)$. De plus
$\vect{CA}(-1\,;0\,;1)$ est orthogonal au vecteur directeur $(1\,;3\,;1)$ de
$(D)$ puisque $-1+0+1 = 0$ : $C$ est bien le projeté orthogonal de $A$, et la
distance vaut $CA = \sqrt2$.

\textbf{3.} On reporte les coordonnées de $(D)$ dans l'équation de $(Q)$ :
$(t+2)-(3t-1)+t = 1$ donne $-t+3 = 1$, soit $t = 2$. Le point cherché est
$I(4\,;5\,;2)$.

\textbf{4.} $\vect{AB}(-3\,;-2\,;0)$ et $\vect{AC}(-5\,;-2\,;-1)$, donc
$\pvect{\vect{AB}}{\vect{AC}} = (2\,;-3\,;-4)$. L'équation du plan est
$2x-3y-4z+9 = 0$. La droite $(\Delta)$ a pour représentation
$x = 4+t$, $y = -1+2t$, $z = 3-t$. Enfin
$\scal{(2;-3;-4)}{(1;2;-1)} = 2-6+4 = 0$ : la direction de $(\Delta)$ est
orthogonale à la normale du plan. Comme $D$ n'appartient pas au plan
— $8+3-12+9 = 8 \neq 0$ —, la droite lui est strictement parallèle.

\textbf{5.} Le vecteur $\vecn(-3\,;-2\,;-6)$ est colinéaire à $(3\,;2\,;6)$ :
les deux plans sont parallèles. $(\mathcal Q)$ a pour équation
$3x+2y+6z-16 = 0$, car $A$ donne $6+4+6 = 16$ ; les constantes diffèrent, donc
le parallélisme est strict. La distance vaut
$\frac{\abs{-33-(-16)}}{\sqrt{9+4+36}} = \frac{49}{7} = 7$. Enfin
$3(-1)+0+6(-5)+33 = 0$ : $H \in (\mathcal P)$, et
$\vect{AH}(-3\,;-2\,;-6)$ est colinéaire à la normale, donc $(AH)$ est
orthogonale à $(\mathcal P)$ ; sa norme vaut $7$.

\textbf{6.} $\pvect{\vecu}{\vecv} = (-6\,;-6\,;6)$, colinéaire à
$(1\,;1\,;-1)$. L'équation est $x+y-z+d = 0$ et $A$ donne $1-2-2+d = 0$, soit
$d = 3$. Le vecteur directeur de $(D)$ est $(-2\,;-2\,;2)$, colinéaire à la
normale : la droite est perpendiculaire au plan. En reportant,
$(-2t+1)+(-2t+3)-(2t-5)+3 = -6t+12 = 0$ donne $t = 2$, d'où
$F(-3\,;-1\,;-1)$.

\textbf{7.} Les directions $(1\,;-2\,;3)$ et $(2\,;1\,;-4)$ ne sont pas
colinéaires. Avec $A(1\,;-1\,;1)$ et $A'(3\,;0\,;1)$, on a
$\vect{AA'}(2\,;1\,;0)$ et $\pvect{\vecu}{\vecu'} = (5\,;10\,;5)$, donc
$\scal{(5;10;5)}{(2;1;0)} = 20 \neq 0$ : les deux droites ne sont pas
coplanaires.

\textbf{8. a)} $\vecn = (5\,;10\,;5)$, colinéaire à $(1\,;2\,;1)$.
\textbf{b)} $(P) : x+2y+z+d = 0$ avec $A(1\,;-1\,;1)$ : $1-2+1+d = 0$, donc
$d = 0$ et $(P) : x+2y+z = 0$.
\textbf{c)} Pour un point de $(D_{1})$ :
$(1+t)+2(-1-2t)+(1+3t) = 1+t-2-4t+1+3t = 0$ : tous les points de $(D_{1})$
appartiennent à $(P)$.

\textbf{9.} $\vect{AB}(1\,;1\,;2)$ et $\vect{AC}(0\,;2\,;2)$ ne sont pas
colinéaires — leur produit vectoriel vaut $(-2\,;-2\,;2)$, non nul — donc les
trois points définissent un plan. Avec la normale $(1\,;1\,;-1)$, l'équation
est $x+y-z+d = 0$ et $A$ donne $-1-1-0+d = 0$, soit $d = 2$ :
$x+y-z+2 = 0$. On vérifie sur $B$ : $0+0-2+2 = 0$, et sur $C$ :
$-1+1-2+2 = 0$.

\textbf{10.} Vraie. Le vecteur $(2\,;3\,;4)$, lu sur l'équation, est normal
au plan, et $(2\,;3\,;4) = 2\left(1\,;\frac32\,;2\right)$ : les deux vecteurs
sont colinéaires, donc $\vecn$ est également normal au plan.
\end{corrige}

% ===========================================================================
\begin{jecherche}[la charpente d'un hangar]
Un hangar a pour base le rectangle $OABC$ avec $O(0\,;0\,;0)$,
$A(12\,;0\,;0)$, $B(12\,;8\,;0)$ et $C(0\,;8\,;0)$, l'unité valant un mètre.
Le faîtage est le segment $[ST]$ avec $S(0\,;4\,;5)$ et $T(12\,;4\,;5)$. Les
deux versants du toit sont les plans $(OAS)$ et $(CBT)$.

\begin{enumerate}[label=\textbf{\arabic*.}, leftmargin=*, itemsep=0.3em]
  \item Déterminer une équation cartésienne de chacun des deux versants.
  \item Vérifier que $T$ appartient au premier versant et que la droite
        $(AT)$ y est contenue.
  \item Calculer la distance du point $B$ au versant $(OAS)$.
  \item Déterminer l'angle que fait le versant $(OAS)$ avec le sol, à un
        dixième de degré près.
  \item Montrer que l'intersection des deux versants est bien la droite
        $(ST)$.
  \item Un câble part de $B$ et suit la direction $\vecu(0\,;-1\,;1)$.
        Rencontre-t-il le faîtage ?
\end{enumerate}
\end{jecherche}

\begin{corrige}[la charpente d'un hangar]
\textbf{1.} Les vecteurs $\vect{OA}(12\,;0\,;0)$ et $\vect{OS}(0\,;4\,;5)$ ont
pour produit vectoriel $(0\,;-60\,;48)$, colinéaire à $(0\,;-5\,;4)$. Le plan passant par
l'origine a donc pour équation $5y-4z = 0$. De même,
$\vect{CB}(12\,;0\,;0)$ et $\vect{CS}(0\,;-4\,;5)$ donnent la normale
$(0\,;5\,;4)$, et $C(0\,;8\,;0)$ conduit à $5y+4z = 40$.

\textbf{2.} $T(12\,;4\,;5)$ : $5 \times 4-4 \times 5 = 0$, donc
$T \in (OAS)$. Comme $A$ y appartient aussi — $0-0 = 0$ —, la droite $(AT)$
est contenue dans le versant.

\textbf{3.} $d\left(B,(OAS)\right) = \dfrac{\abs{5 \times 8-0}}{\sqrt{25+16}}
= \dfrac{40}{\sqrt{41}} \approx 6{,}25$~mètres.

\textbf{4.} Le sol a pour normale $(0\,;0\,;1)$ et le versant $(0\,;-5\,;4)$.
L'angle $\theta$ des deux plans vérifie
$\cos\theta = \dfrac{4}{\sqrt{41}} \approx 0{,}6247$, soit
$\theta \approx 51{,}3^{\circ}$ — ce que confirme la pente, qui monte de
$5$~mètres pour $4$~mètres d'avancée.

\textbf{5.} Un point appartient aux deux versants si $5y-4z = 0$ et
$5y+4z = 40$. En additionnant, $10y = 40$ donc $y = 4$, puis $z = 5$ ;
l'abscisse reste libre. C'est exactement la droite passant par $S$ et $T$.

\textbf{6.} Le câble a pour représentation $x = 12$, $y = 8-s$, $z = s$. Le
faîtage est l'ensemble des points $(x\,;4\,;5)$. L'égalité imposerait
$s = 5$ et $8-s = 4$, c'est-à-dire $s = 3$ : impossible. Le câble ne
rencontre donc pas le faîtage, et les deux droites ne sont pas coplanaires.
\end{corrige}

% ===========================================================================
\begin{jemeteste}
Répondre par vrai ou faux, en justifiant en une ligne.

\begin{enumerate}[label=\textbf{\arabic*.}, leftmargin=*, itemsep=0.28em]
  \item Deux droites de l'espace qui ne sont pas parallèles sont sécantes.
  \item Le produit vectoriel de deux vecteurs colinéaires est nul.
  \item Si $\scal{\vecn}{\vecu} = 0$, la droite dirigée par $\vecu$ est
        contenue dans le plan de normale $\vecn$.
  \item $\pvect{\vecu}{\vecv} = \pvect{\vecv}{\vecu}$.
  \item Deux plans dont les normales sont colinéaires sont parallèles.
  \item L'équation $x+y+z = 0$ définit un plan passant par l'origine.
\end{enumerate}
\end{jemeteste}

\begin{corrige}[je me teste]
\textbf{1.} Faux — elles peuvent être non coplanaires.
\textbf{2.} Vrai.
\textbf{3.} Faux — elle est parallèle au plan, mais pas nécessairement
contenue dedans ; il faut vérifier qu'un point lui appartient.
\textbf{4.} Faux — les deux produits sont opposés.
\textbf{5.} Vrai — confondus ou strictement parallèles.
\textbf{6.} Vrai — les coordonnées de l'origine vérifient l'équation.
\end{corrige}

% ===========================================================================
\begin{commeaubac}
\bmtsource{Baccalauréat, série S, session 2026 — exercice 2}
Dans un espace orienté muni d'un repère orthonormé direct, on considère le
plan $(P)$ contenant le point $A(1\,;-2\,;2)$ et dirigé par les vecteurs
$\vecu(-1\,;2\,;1)$ et $\vecv(-4\,;2\,;-2)$, ainsi que la droite $(D)$ de
représentation paramétrique
\[
  (D) : \begin{cases} x = -2t+1 \\ y = -2t+3 \\ z = 2t-5 \end{cases}
  \qquad t \in \R .
\]

\begin{enumerate}[label=\textbf{\arabic*.}, leftmargin=*, itemsep=0.25em]
  \item Montrer que $(P)$ a pour équation cartésienne $x+y-z+3 = 0$.
        \hfill\textup{(0,5 pt)}
  \item Vérifier que $(P)$ et $(D)$ sont perpendiculaires.
        \hfill\textup{(0,5 pt)}
  \item Déterminer les coordonnées du point $F$, intersection de $(P)$ et
        $(D)$. \hfill\textup{(0,5 pt)}
  \item Soit $(D')$ la droite passant par $A$ et perpendiculaire à $(D)$.
        Déterminer la distance minimale entre $(D)$ et $(D')$.
        \hfill\textup{(0,5 pt)}
\end{enumerate}
\end{commeaubac}

\begin{corrige}[comme au baccalauréat]
\textbf{1.} Le produit vectoriel $\pvect{\vecu}{\vecv}$ a pour coordonnées
\[
  \bigl(2 \times (-2)-1 \times 2\,;\ 1 \times (-4)-(-1) \times (-2)\,;\
  (-1) \times 2-2 \times (-4)\bigr),
\]
soit $(-6\,;-6\,;6)$, colinéaire à $(1\,;1\,;-1)$. L'équation est donc $x+y-z+d = 0$, et $A$ donne
$1-2-2+d = 0$, soit $d = 3$.

\textbf{2.} Le vecteur directeur de $(D)$ est $(-2\,;-2\,;2)$, colinéaire au
vecteur normal $(1\,;1\,;-1)$ : la droite est perpendiculaire au plan.

\textbf{3.} En reportant : $-6t+12 = 0$, donc $t = 2$ et
$F(-3\,;-1\,;-1)$.

\textbf{4.} $(D')$ est perpendiculaire à $(D)$ et passe par $A$, qui
appartient à $(P)$ ; or $(D)$ est perpendiculaire à $(P)$, donc toute droite
de $(P)$ passant par $A$ convient. La distance minimale entre $(D)$ et
$(D')$ est donc la distance de $A$ au point $F$ :
\[
  AF = \sqrt{(-3-1)^{2}+(-1+2)^{2}+(-1-2)^{2}} = \sqrt{26}.
\]
\end{corrige}

% ===========================================================================
\begin{concours}
\bmtsource{D'après un concours d'entrée en école d'ingénieurs}
Dans un repère orthonormé, on considère les points $A(1\,;0\,;0)$,
$B(0\,;1\,;0)$, $C(0\,;0\,;1)$ et $D(1\,;1\,;1)$.

\begin{enumerate}[label=\textbf{\arabic*.}, leftmargin=*, itemsep=0.25em]
  \item Montrer que $ABC$ est équilatéral et calculer son aire.
  \item Déterminer une équation du plan $(ABC)$ et la distance de $D$ à ce
        plan.
  \item En déduire le volume du tétraèdre $ABCD$.
\end{enumerate}
\end{concours}

\begin{corrige}[vers le concours]
\textbf{1.} $AB = AC = BC = \sqrt2$ : le triangle est équilatéral. Son aire
vaut $\frac{\sqrt3}{4}\left(\sqrt2\right)^{2} = \frac{\sqrt3}{2}$ ; on la
retrouve par $\frac12\norme{\pvect{\vect{AB}}{\vect{AC}}}$, avec
$\pvect{\vect{AB}}{\vect{AC}} = (1\,;1\,;1)$ de norme $\sqrt3$.

\textbf{2.} Le vecteur $(1\,;1\,;1)$ est normal, et $A$ donne
$(ABC) : x+y+z-1 = 0$. La distance de $D(1\,;1\,;1)$ vaut
$\frac{\abs{1+1+1-1}}{\sqrt3} = \frac{2}{\sqrt3} = \frac{2\sqrt3}{3}$.

\textbf{3.} $V = \frac13 \times \text{aire} \times \text{hauteur}
= \frac13 \times \frac{\sqrt3}{2} \times \frac{2\sqrt3}{3} = \frac13$.
\end{corrige}

% ===========================================================================
\begin{notepourlaclasse}
Le chapitre est mécanique, et c'est une bonne nouvelle : bien travaillé, il
rapporte deux points en vingt minutes. Consacrez la première séance au seul
produit vectoriel, avec la disposition en croix et la vérification
systématique par produit scalaire. Un élève qui calcule juste ce produit
traite l'exercice entier.

Le point conceptuel unique est celui des droites non coplanaires. Faites-le
sentir avec un objet : deux crayons tenus en l'air, qui ne se touchent pas et
ne sont pas parallèles. Sans cette image, les élèves concluent
systématiquement « non parallèles donc sécantes », et perdent la question.

Le vocabulaire mérite d'être fixé au tableau : parallèle, strictement
parallèle, confondu, perpendiculaire, orthogonal. Les sujets l'emploient avec
précision, et une réponse juste rédigée avec le mauvais mot est comptée
fausse.
\end{notepourlaclasse}

% =========================================================================
\begin{devoir}[2 heures]
Trois exercices indépendants, notés sur $20$. Le devoir est cumulatif : il porte sur ce chapitre et sur ceux qui le précèdent. Les énoncés sont repris de sessions réelles et assemblés ici en épreuve ; leur provenance est indiquée à chaque fois.

\exods{1}{Baccalauréat, Cameroun, séries D et TI, session 2024 — exercice 1}{5 points}
Le plan complexe est muni d'un repère orthonormé. On donne $A$, $B$ et $C$
d'affixes respectives $Z_{A} = 2-i$, $Z_{B} = 3-2i$ et $Z_{C} = 1-2i$.
\begin{enumerate}[label=\textbf{\alph*)}, leftmargin=*, itemsep=0.15em]
  \item Déterminer l'expression complexe de la similitude directe $s$ de
        centre $C$ transformant $A$ en $B$.
  \item Préciser les éléments caractéristiques de $s$.
  \item Quelle est l'image de la droite $(AC)$ par $s$ ?
\end{enumerate}

\exods{2}{Baccalauréat, Côte d'Ivoire, série D, session 2023 — exercice 4}{7 points}
On donne $\Omega(1+i)$, $A(1)$ et $B\!\left(\frac32+\frac i2\right)$. Étudier
la similitude directe $S$ de centre $\Omega$ transformant $A$ en $B$ :
déterminer son écriture complexe, puis son rapport et son angle.

\exods{3}{Baccalauréat, Cameroun, série C, session 2023 — exercice 3}{8 points}
L'espace est muni d'un repère orthonormé direct. Soient $A(-1\,;-1\,;0)$,
$B(0\,;0\,;2)$ et $C(-1\,;1\,;2)$. Montrer que $A$, $B$ et $C$ définissent un
plan, puis en donner une équation cartésienne.

\end{devoir}

\begin{corrige}[devoir surveillé]
\textbf{Exercice 1. a)} $a = \dfrac{Z_{B}-Z_{C}}{Z_{A}-Z_{C}} = \dfrac{2}{1+i} = 1-i$,
donc $z' = (1-i)(z-1+2i)+1-2i$.
\textbf{b)} Centre $C(1-2i)$, rapport $\abs{1-i} = \sqrt2$, angle
$-\frac{\pi}{4}$.
\textbf{c)} Une similitude transforme une droite en droite. Comme
$s(C) = C$ et $s(A) = B$, l'image de $(AC)$ est la droite $(BC)$.

\textbf{Exercice 2.} $a = \dfrac{z_{B}-z_{\Omega}}{z_{A}-z_{\Omega}}
= \dfrac{\frac32+\frac i2-1-i}{1-1-i}
= \dfrac{\frac12-\frac i2}{-i} = \dfrac{1+i}{2}$. L'écriture complexe est
$z' = \frac{1+i}{2}\left(z-1-i\right)+1+i$, le rapport vaut
$\left|\frac{1+i}{2}\right| = \frac{\sqrt2}{2}$ et l'angle
$\frac{\pi}{4}$.

\textbf{Exercice 3.} $\vect{AB}(1\,;1\,;2)$ et $\vect{AC}(0\,;2\,;2)$ ne sont pas
colinéaires — leur produit vectoriel vaut $(-2\,;-2\,;2)$, non nul — donc les
trois points définissent un plan. Avec la normale $(1\,;1\,;-1)$, l'équation
est $x+y-z+d = 0$ et $A$ donne $-1-1-0+d = 0$, soit $d = 2$ :
$x+y-z+2 = 0$. On vérifie sur $B$ : $0+0-2+2 = 0$, et sur $C$ :
$-1+1-2+2 = 0$.

\end{corrige}
